% Reconstructed text from V5 PDF page 48; non-native source
\par\smallskip\noindent\rule{\textwidth}{0.4pt}\par{\sffamily\footnotesize Original V5 page 48 -- recovered text}\par\begin{flushleft}\footnotesize\ttfamily
\par\smallskip
\par\smallskip
\noindent Diagnostics and Serving Infrastructure for MAS 48\par
\par\smallskip
\par\smallskip
\par\smallskip
\noindent The subsampling paradox reveals heuristic limitations. The retention ablation results (Sec-\par
\noindent tion 7.3) demonstrate that heuristic attributors can improve when given less data — a finding\par
\noindent that challenges the intuition that more evidence is always better for diagnosis. The mechanism\par
\noindent is specific: the causal-graph heuristic marks an error as “recovered” if the same agent later pro-\par
\noindent duces a normal event, causing it to skip the true root cause on full traces where normal post-error\par
\noindent events are abundant. Subsampling strips these misleading normal events, making the decisive\par
\noindent error visible. This finding has two implications for cross-layer design: first, retention policies\par
\noindent that prioritize recency (keeping the most recent events) perform worst at low budgets because\par
\noindent they systematically discard early root causes. Second, the optimal retention strategy depends\par
\noindent on the attribution method — a finding that argues against one-size-fits-all retention policies and\par
\noindent for the retention-aware diagnostic design advocated in our validation protocol.\par
\par\smallskip
\noindent 9.2 Implications for System Design\par
\noindent The patterns above suggest three concrete architectural recommendations:\par
\noindent  1. Decouple diagnostic storage from serving memory. Since 10 of 12 failure-localization\par
\noindent  methods operate on logged traces rather than live serving state, the simplest way to avoid\par
\noindent  Tier C evidence-loss events is to ensure traces are written to a durable store before the\par
\noindent  corresponding KV cache is eligible for eviction. This is already the default architecture\par
\noindent  in most surveyed systems, which is why the cross-layer conflict is rare in practice.\par
\noindent  2. Expose serving-side events to diagnostic consumers. The observability gap (no OTel\par
\noindent  attributes for KV cache residency, eviction events, or memory pressure) prevents diag-\par
\noindent  nostic systems from detecting when their evidence might be at risk. Extending the OTel\par
\noindent  GenAI conventions with serving-side attributes would enable diagnostic systems to re-\par
\noindent  quest trace retention proactively when they detect that a future attribution may be needed.\par
\noindent  3. Design retention policies with attribution method awareness. The subsampling para-\par
\noindent  dox shows that the same retention strategy can help one attribution method and hurt an-\par
\noindent  other. Future serving systems that incorporate diagnostic retention should parameterize\par
\noindent  the retention policy on the expected attribution method, or provide a general-purpose\par
\noindent  noise-filtering retention strategy that benefits multiple methods.\par
\par\smallskip
\noindent 9.3 Limitations\par
\noindent This review has several limitations:\par
\noindent  1. Curated, not systematic. The corpus was assembled through iterative exploration and\par
\noindent  citation chaining rather than a pre-registered protocol with reproducible database queries.\par
\noindent  The reported counts are a curated landscape, not an exhaustive census.\par
\noindent  2. Evidence depth variation. Of the 60 entries, 56 received full-text verification (T1a or\par
\noindent  T1b), comprising 44 systems and all 12 benchmarks; the remaining 4 entries are discussed\par
\noindent  at the T2 text-mention level. Cross-system comparisons involving T2 systems should be\par
\noindent  interpreted with this asymmetry in mind.\par
\noindent  3. Retention ablation scope. The retention-ablation experiment uses real TraceElephant\par
\noindent  traces (220 failures, 5 task suites from 3 MAS frameworks) but tests only two heuristic at-\par
\noindent  tribution methods. LLM-judge approaches (CHIEF, TrajAudit, AgentRx) were not tested\par
\noindent  and may not exhibit the subsampling paradox, since they do not rely on the recovery-\par
\noindent  detection heuristic that drives the effect. Step-level attribution was not tested.\par
\noindent  4. Rapidly evolving landscape. Multi-agent runtime research is moving faster than review\par
\noindent  cycles. Several systems in the corpus are preprints without peer review. The search cutoff\par
\end{flushleft}
